% =============================================================================
% Aula 06 — A classe do meio · Empréstimo na biblioteca · Venda no StockSales
% Beamer + tema Madrid | Curso Entra21 — Spring Framework
% Compilar: tectonic -X compile slides-aula-06.tex
% =============================================================================
\documentclass[aspectratio=169,11pt]{beamer}

\usetheme{Madrid}
\usecolortheme{default}

\usepackage[T1]{fontenc}
\usepackage[utf8]{inputenc}
\usepackage[brazil]{babel}
\usepackage{booktabs}
\usepackage{graphicx}
\usepackage{hyperref}
\usepackage{listings}
\usepackage{xcolor}
\usepackage{tikz}
\usetikzlibrary{arrows.meta, positioning, shapes.geometric}

\definecolor{codebg}{RGB}{245,245,245}
\definecolor{codegreen}{RGB}{40,120,40}
\definecolor{codeblue}{RGB}{30,70,160}
\definecolor{codered}{RGB}{180,50,50}

\lstset{
  basicstyle=\ttfamily\small,
  backgroundcolor=\color{codebg},
  frame=single,
  rulecolor=\color{black!20},
  breaklines=true,
  columns=fullflexible,
  keepspaces=true,
  showstringspaces=false,
  keywordstyle=\color{codeblue}\bfseries,
  commentstyle=\color{codegreen},
  stringstyle=\color{codered},
  tabsize=2,
  literate={á}{{\'a}}1 {é}{{\'e}}1 {í}{{\'i}}1 {ó}{{\'o}}1 {ú}{{\'u}}1
           {Á}{{\'A}}1 {É}{{\'E}}1 {Í}{{\'I}}1 {Ó}{{\'O}}1 {Ú}{{\'U}}1
           {ã}{{\~a}}1 {õ}{{\~o}}1 {Ã}{{\~A}}1 {Õ}{{\~O}}1
           {ç}{{\c{c}}}1 {Ç}{{\c{C}}}1
           {ê}{{\^e}}1 {Ê}{{\^E}}1 {â}{{\^a}}1 {Â}{{\^A}}1
           {à}{{\`a}}1 {À}{{\`A}}1
}

\newcommand{\onde}[1]{%
  \begin{block}{Onde estamos}
    #1
  \end{block}
}

\title{Aula 06 --- Spring Framework}
\subtitle{A classe do meio · Empréstimo · Venda}
\author{Mauricio Duailibi Neto}
\institute{Turma Java Entra21}
\date{}

\begin{document}

\begin{frame}
  \titlepage
\end{frame}

\begin{frame}{Como será a aula de hoje}
  \begin{block}{Parte 1 --- 18h30 às 20h}
    \begin{itemize}
      \item Revisão da Aula 05
      \item Quando 1{:}N não chega
      \item A classe do meio
      \item Projeto \textbf{biblioteca} --- fazemos juntos
    \end{itemize}
  \end{block}

  \vspace{0.4em}
  \begin{block}{Parte 2 --- 20h20 às 22h}
    \begin{itemize}
      \item Voltamos ao \textbf{StockSales}
      \item Venda de um produto
      \item Venda com vários produtos
    \end{itemize}
  \end{block}
\end{frame}

\begin{frame}{Dois projetos nesta aula}
  Não vamos misturar os códigos.

  \vspace{0.6em}
  \begin{description}
    \item[biblioteca] Primeira parte da aula.
      Livro e leitor já existem. Nós criamos o empréstimo.
    \item[StockSales] O sistema da loja.
      Cliente e produto já existem. Só depois do intervalo.
  \end{description}

  \vspace{0.6em}
  O JavaScript das telas já está pronto.\\
  Hoje a gente escreve o Java.
\end{frame}

\begin{frame}{Mapa da aula}
  \tableofcontents
\end{frame}

% #############################################################################
\section{Revisão da Aula 05}

\begin{frame}{O que já sabemos}
  Na aula passada um cadastro passou a apontar para outro.

  \vspace{0.5em}
  \begin{itemize}
    \item laboratorio: o material guarda \texttt{categoriaId}
    \item Uma categoria, vários materiais --- isso é 1{:}N
    \item As listas ficam na classe \texttt{Dados}
    \item StockSales: produto e cliente, cada um na sua lista
  \end{itemize}

  \vspace{0.5em}
  Cliente e produto ainda não se encontram.\\
  Ainda não existe venda.
\end{frame}

\begin{frame}{Como gravamos o 1{:}N}
  Não copiamos a categoria inteira dentro do material.

  \vspace{0.5em}
  Guardamos só o número:

  \vspace{0.4em}
  \begin{center}
  \begin{tabular}{lll}
    \toprule
    Material & quantidade & categoriaId \\
    \midrule
    Caneta & 10 & 1 \\
    Caderno & 5 & 1 \\
    Mouse & 3 & 2 \\
    \bottomrule
  \end{tabular}
  \end{center}

  \vspace{0.5em}
  A categoria 1 é Papelaria.\\
  Se o nome da categoria mudar, o material continua ligado ao id 1.
\end{frame}

\begin{frame}{Os quatro verbos continuam iguais}
  \begin{center}
  \begin{tabular}{lll}
    \toprule
    Ação & Verbo HTTP & O que faz \\
    \midrule
    Listar / ler & GET & Lê, não muda nada \\
    Cadastrar & POST & Cria um item novo \\
    Editar & PUT & Troca os dados de um item \\
    Excluir & DELETE & Apaga um item \\
    \bottomrule
  \end{tabular}
  \end{center}

  \vspace{0.7em}
  Hoje o POST ganha um trabalho extra:\\
  achar os dois cadastros e conferir a quantidade.
\end{frame}

% #############################################################################
\section{A classe do meio}

\begin{frame}{O problema de hoje}
  Na biblioteca, Maria quer levar Dom Casmurro.

  \vspace{0.5em}
  Não dá para colocar o livro dentro do leitor.\\
  O leitor levaria um livro só, para sempre.

  \vspace{0.5em}
  Não dá para colocar o leitor dentro do livro.\\
  O livro serviria para uma pessoa só.

  \vspace{0.5em}
  \begin{block}{Pergunta de hoje}
    Onde gravamos o encontro dos dois,\\
    com a quantidade que saiu da prateleira?
  \end{block}
\end{frame}

\begin{frame}{Por que um id só não resolve}
  No 1{:}N, o material aponta para \textbf{uma} categoria.

  \vspace{0.5em}
  O empréstimo aponta para \textbf{dois} lados:
  \begin{itemize}
    \item quem leva --- o leitor
    \item o que leva --- o livro
  \end{itemize}

  \vspace{0.5em}
  E ainda precisa de coisa própria:
  \begin{itemize}
    \item um id
    \item a quantidade
  \end{itemize}

  \vspace{0.4em}
  Isso não cabe em \texttt{leitorId} dentro do livro.\\
  Nem em \texttt{livroId} dentro do leitor.
\end{frame}

\begin{frame}{A classe do meio}
  Criamos um cadastro novo.\\
  Cada linha é um encontro.

  \vspace{0.4em}
  \begin{center}
  \begin{tabular}{llll}
    \toprule
    id & leitorId & livroId & quantidade \\
    \midrule
    1 & 1 Maria & 1 Dom Casmurro & 1 \\
    2 & 3 Helena & 3 Capitaes da Areia & 2 \\
    \bottomrule
  \end{tabular}
  \end{center}

  \vspace{0.5em}
  Maria continua na lista de leitores.\\
  O livro continua na lista de livros.\\
  O empréstimo só guarda os ids e a quantidade.
\end{frame}

\begin{frame}{O desenho}
  \begin{center}
  \begin{tikzpicture}[
      box/.style={draw, rounded corners, align=center,
                  minimum width=2.6cm, minimum height=0.9cm,
                  font=\small},
      seta/.style={-{Latex}, thick}
    ]
    \node[box, fill=blue!15] (lei) {Leitor};
    \node[box, fill=yellow!25, right=2.0cm of lei] (emp) {Emprestimo};
    \node[box, fill=green!20, right=2.0cm of emp] (liv) {Livro};

    \draw[seta] (lei) -- (emp);
    \draw[seta] (liv) -- (emp);
  \end{tikzpicture}
  \end{center}

  \vspace{0.6em}
  O amarelo é a classe do meio.\\
  Sem ela, leitor e livro não se encontram.

  \vspace{0.4em}
  Um leitor tem vários empréstimos.\\
  Um livro aparece em vários empréstimos.\\
  Cada empréstimo, neste primeiro passo, aponta para \textbf{um} livro.
\end{frame}

\begin{frame}{Dois passos na aula}
  \begin{center}
  \begin{tabular}{ll}
    \toprule
    Passo & O que a classe do meio guarda \\
    \midrule
    1 & Um leitor e \textbf{um} livro, com quantidade \\
    2 & Um leitor e uma \textbf{lista} de livros \\
    \bottomrule
  \end{tabular}
  \end{center}

  \vspace{0.7em}
  Primeiro a gente faz o passo 1 funcionar na tela.\\
  Só depois a lista.

  \vspace{0.5em}
  No StockSales é a mesma ordem:\\
  um produto, depois vários produtos.
\end{frame}

\begin{frame}{O que não vamos usar hoje}
  \begin{itemize}
    \item Banco de dados
    \item Anotações do tipo \texttt{@ManyToMany}
    \item JPA
    \item Data de empréstimo
  \end{itemize}

  \vspace{0.5em}
  Continuamos com lista na memória e um \texttt{int}\\
  guardando o id do outro cadastro.

  \vspace{0.5em}
  O raciocínio é o mesmo que o banco vai usar depois.\\
  Só a ferramenta muda.
\end{frame}

% #############################################################################
\section{Projeto biblioteca}

\begin{frame}{O projeto biblioteca}
  Projeto novo. Não é o laboratorio.

  \vspace{0.5em}
  \begin{block}{O que já está pronto}
    Livro e leitor com GET, POST, PUT e DELETE.\\
    Telas com formulário. Axios já chama a API de empréstimo.
  \end{block}

  \vspace{0.5em}
  Se a API ainda não existe, a tela avisa.\\
  Quando o Java aparecer, a lista responde.
\end{frame}

\begin{frame}{Passo 1 --- Subir o Spring}
  \onde{Projeto \textbf{biblioteca}}

  No terminal, dentro da pasta \texttt{biblioteca}:

  \vspace{0.4em}
  \texttt{./mvnw spring-boot:run}\\
  {\small No Windows: \texttt{.\textbackslash mvnw.cmd spring-boot:run}}

  \vspace{0.5em}
  Abram no navegador: \texttt{http://localhost:8080}

  \vspace{0.5em}
  \begin{alertblock}{Um Spring por vez}
    Se outro projeto ainda estiver na porta 8080,
    parem o outro antes.
  \end{alertblock}
\end{frame}

\begin{frame}{O que vocês vão ver na tela}
  \begin{itemize}
    \item Livros: Dom Casmurro, O Cortico, Capitaes da Areia, A Hora da Estrela
    \item Leitores: Maria, Jose, Helena
    \item Novo empréstimo: dois selects e a quantidade
    \item Histórico: lista vazia, por enquanto
  \end{itemize}

  \vspace{0.5em}
  Cadastrar livro e leitor já funciona.\\
  Confirmar empréstimo ainda não.\\
  Esse Java a gente escreve agora.
\end{frame}

\begin{frame}[fragile]{Onde cada arquivo fica}
\begin{lstlisting}[basicstyle=\ttfamily\footnotesize]
biblioteca/
  src/main/java/com/entra21/biblioteca/
    BibliotecaApplication.java
    Dados.java
    model/Livro.java
    model/Leitor.java
    model/Emprestimo.java          novo
    controller/LivroController.java
    controller/LeitorController.java
    controller/EmprestimoController.java  novo
  src/main/resources/static/
    emprestimos-novo.html
    emprestimos.html
\end{lstlisting}
\end{frame}

\begin{frame}[fragile]{Emprestimos em Dados}
  \onde{Projeto \textbf{biblioteca}\\
        Arquivo \texttt{Dados.java}, acrescentar}

\begin{lstlisting}[language=Java,basicstyle=\ttfamily\footnotesize]
public static ArrayList<Emprestimo> emprestimos =
    new ArrayList<>();
public static int proximoIdEmprestimo = 1;
\end{lstlisting}

  \vspace{0.4em}
  Import de \texttt{Emprestimo} no topo do arquivo.

  \vspace{0.4em}
  Sem isso, o Controller não tem onde guardar.
\end{frame}

\begin{frame}[fragile]{Classe Emprestimo}
  \onde{Arquivo novo:\\
        \texttt{src/main/java/com/entra21/biblioteca/model/Emprestimo.java}}

\begin{lstlisting}[language=Java,basicstyle=\ttfamily\footnotesize]
public class Emprestimo {
  private int id;
  private int leitorId;
  private int livroId;
  private int quantidade;

  public Emprestimo() {
  }
}
\end{lstlisting}

  Construtor vazio e os gets e sets dos quatro atributos.\\
  O vazio é obrigatório no POST.
\end{frame}

\begin{frame}[fragile]{O JSON que a tela já manda}
  Vocês não mexem no JavaScript.\\
  A página já envia isto:

  \vspace{0.4em}
\begin{lstlisting}[basicstyle=\ttfamily\footnotesize]
{"leitorId": 1, "livroId": 2, "quantidade": 1}
\end{lstlisting}

  \vspace{0.5em}
  O servidor faz o resto:
  \begin{enumerate}
    \item Achar o leitor
    \item Achar o livro
    \item Se a quantidade do livro for menor, recusar
    \item Diminuir a quantidade do livro
    \item Guardar o empréstimo
  \end{enumerate}
\end{frame}

\begin{frame}[fragile]{Uma função só para achar}
  \onde{\texttt{controller/EmprestimoController.java}, arquivo novo}

  Em vez de um \texttt{for} gigante dentro do POST,\\
  uma função pequena devolve o livro ou \texttt{null}.

\begin{lstlisting}[language=Java,basicstyle=\ttfamily\scriptsize]
public Livro acharLivro(int id) {
  for (int i = 0; i < Dados.livros.size(); i++) {
    if (Dados.livros.get(i).getId() == id) {
      return Dados.livros.get(i);
    }
  }
  return null;
}
\end{lstlisting}

  Façam a mesma ideia para \texttt{acharLeitor}.
\end{frame}

\begin{frame}[fragile]{GET dos empréstimos}
  \onde{Ainda em \texttt{EmprestimoController.java}}

\begin{lstlisting}[language=Java,basicstyle=\ttfamily\footnotesize]
@RestController
public class EmprestimoController {

  @GetMapping("/api/emprestimos")
  public ArrayList<Emprestimo> listar() {
    return Dados.emprestimos;
  }
}
\end{lstlisting}

  \vspace{0.4em}
  Testem no navegador:\\
  \texttt{http://localhost:8080/api/emprestimos}

  \vspace{0.3em}
  Deve aparecer \texttt{[]}. Lista vazia, mas a rota existe.
\end{frame}

\begin{frame}[fragile]{POST --- achar leitor e livro}
  \onde{Ainda no método \texttt{cadastrar}}

\begin{lstlisting}[language=Java,basicstyle=\ttfamily\scriptsize]
@PostMapping("/api/emprestimos")
public Emprestimo cadastrar(@RequestBody Emprestimo pedido) {
  Leitor leitor = acharLeitor(pedido.getLeitorId());
  if (leitor == null) {
    throw new ResponseStatusException(
        HttpStatus.BAD_REQUEST, "Leitor nao encontrado");
  }

  Livro livro = acharLivro(pedido.getLivroId());
  if (livro == null) {
    throw new ResponseStatusException(
        HttpStatus.BAD_REQUEST, "Livro nao encontrado");
  }
\end{lstlisting}

  O método ainda não acabou.
\end{frame}

\begin{frame}[fragile]{POST --- quantidade e gravação}
  \onde{Ainda no mesmo método \texttt{cadastrar}}

\begin{lstlisting}[language=Java,basicstyle=\ttfamily\scriptsize]
  if (livro.getQuantidade() < pedido.getQuantidade()) {
    throw new ResponseStatusException(
        HttpStatus.BAD_REQUEST, "Quantidade insuficiente");
  }

  pedido.setId(Dados.proximoIdEmprestimo);
  Dados.proximoIdEmprestimo = Dados.proximoIdEmprestimo + 1;
  livro.setQuantidade(
      livro.getQuantidade() - pedido.getQuantidade());
  Dados.emprestimos.add(pedido);
  return pedido;
}
\end{lstlisting}
\end{frame}

\begin{frame}{Imports deste Controller}
  \begin{itemize}
    \item \texttt{Dados}, \texttt{Emprestimo}, \texttt{Leitor}, \texttt{Livro}
    \item \texttt{HttpStatus}
    \item \texttt{ResponseStatusException}
    \item \texttt{ArrayList}
    \item As anotações \texttt{GetMapping}, \texttt{PostMapping},
          \texttt{RequestBody}, \texttt{RestController}
  \end{itemize}

  \vspace{0.5em}
  Se faltar um import, o VS Code sublinha em vermelho.\\
  Aceitem o import que ele oferece.
\end{frame}

\begin{frame}{Conferindo o passo 1}
  \begin{itemize}
    \item [ ] \texttt{/api/emprestimos} devolve \texttt{[]}
    \item [ ] Maria leva 1 Dom Casmurro --- a tela diz registrado
    \item [ ] Em Livros, a quantidade do titulo caiu
    \item [ ] O histórico mostra o nome, não só o id
    \item [ ] Pedir mais exemplares do que existe --- a tela recusa
  \end{itemize}

  \vspace{0.5em}
  Se a quantidade não cair, o \texttt{setQuantidade} não rodou.\\
  Se o histórico continuar vazio, o GET não devolveu a lista.
\end{frame}

\begin{frame}{Erros que mais aparecem}
  \begin{itemize}
    \item 404 no POST --- o Controller não compilou ou o Spring
          não reiniciou
    \item Lista de livros some --- o seed não está no
          \texttt{if (Dados.livros.isEmpty())}
    \item Quantidade não diminui --- acharam o livro e não
          gravaram de volta no objeto
    \item Porta 8080 ocupada
  \end{itemize}
\end{frame}

% #############################################################################
\section{Vários livros}

\begin{frame}{Um empréstimo, vários livros}
  Na vida real a pessoa sai com mais de um título.

  \vspace{0.5em}
  Um campo \texttt{livroId} não basta.\\
  Precisamos de uma lista dentro do empréstimo.

  \vspace{0.5em}
  Cada linha da lista é um \textbf{item}:\\
  este livro, esta quantidade.

  \vspace{0.5em}
  O POST de um livro \textbf{continua no ar}.\\
  Vamos acrescentar outro endereço.
\end{frame}

\begin{frame}{O desenho com lista}
  \begin{center}
  \begin{tikzpicture}[
      box/.style={draw, rounded corners, align=center,
                  minimum width=2.2cm, minimum height=0.8cm,
                  font=\small},
      meio/.style={draw, rounded corners, align=center,
                   minimum width=1.8cm, minimum height=0.55cm,
                   font=\scriptsize},
      seta/.style={-{Latex}, thick}
    ]
    \node[box, fill=blue!15] (lei) {Helena};
    \node[box, fill=yellow!25, right=1.6cm of lei] (emp) {Emprestimo};
    \node[meio, fill=orange!25, right=1.4cm of emp] (i1) {item 1};
    \node[meio, fill=orange!25, below=0.25cm of i1] (i2) {item 2};
    \node[box, fill=green!20, right=1.4cm of i1] (l1) {Dom Casmurro};
    \node[box, fill=green!20, below=0.25cm of l1] (l2) {O Cortico};

    \draw[seta] (lei) -- (emp);
    \draw[seta] (emp) -- (i1);
    \draw[seta] (emp) -- (i2);
    \draw[seta] (i1) -- (l1);
    \draw[seta] (i2) -- (l2);
  \end{tikzpicture}
  \end{center}

  \vspace{0.3em}
  Um empréstimo. Dois itens. Dois livros.
\end{frame}

\begin{frame}[fragile]{Classe ItemEmprestimo}
  \onde{Arquivo novo:\\
        \texttt{model/ItemEmprestimo.java}}

\begin{lstlisting}[language=Java,basicstyle=\ttfamily\footnotesize]
public class ItemEmprestimo {
  private int livroId;
  private int quantidade;

  public ItemEmprestimo() {
  }
}
\end{lstlisting}

  Gets e sets dos dois atributos.\\
  Esta classe não precisa de id próprio nesta aula.
\end{frame}

\begin{frame}[fragile]{Lista de itens no Emprestimo}
  \onde{\texttt{model/Emprestimo.java}, acrescentar}

\begin{lstlisting}[language=Java,basicstyle=\ttfamily\footnotesize]
private ArrayList<ItemEmprestimo> itens;

public ArrayList<ItemEmprestimo> getItens() {
  return itens;
}

public void setItens(ArrayList<ItemEmprestimo> itens) {
  this.itens = itens;
}
\end{lstlisting}

  \vspace{0.4em}
  Os campos \texttt{livroId} e \texttt{quantidade} ficam.\\
  O POST de um livro ainda usa eles.
\end{frame}

\begin{frame}[fragile]{O JSON da tela de vários livros}
  A página \texttt{emprestimos-varios.html} já manda:

  \vspace{0.3em}
\begin{lstlisting}[basicstyle=\ttfamily\footnotesize]
{
  "leitorId": 3,
  "itens": [
    {"livroId": 1, "quantidade": 1},
    {"livroId": 2, "quantidade": 1}
  ]
}
\end{lstlisting}

  \vspace{0.4em}
  Endereço: POST \texttt{/api/emprestimos/itens}
\end{frame}

\begin{frame}{Dois fors, nessa ordem}
  Primeiro conferimos \textbf{tudo}.\\
  Só depois baixamos quantidade.

  \vspace{0.5em}
  Se o segundo livro não tiver exemplar,\\
  o primeiro também não sai.

  \vspace{0.5em}
  \begin{enumerate}
    \item \texttt{for} nos itens: achar o livro, conferir quantidade
    \item Se algum falhar, recusar o empréstimo inteiro
    \item Outro \texttt{for}: diminuir a quantidade de cada livro
    \item Gravar o empréstimo
  \end{enumerate}
\end{frame}

\begin{frame}[fragile]{POST de vários --- conferir}
  \onde{\texttt{EmprestimoController.java}, método novo}

\begin{lstlisting}[language=Java,basicstyle=\ttfamily\scriptsize]
@PostMapping("/api/emprestimos/itens")
public Emprestimo cadastrarVarios(@RequestBody Emprestimo pedido) {
  Leitor leitor = acharLeitor(pedido.getLeitorId());
  if (leitor == null) {
    throw new ResponseStatusException(
        HttpStatus.BAD_REQUEST, "Leitor nao encontrado");
  }

  for (int i = 0; i < pedido.getItens().size(); i++) {
    ItemEmprestimo item = pedido.getItens().get(i);
    Livro livro = acharLivro(item.getLivroId());
    if (livro == null) {
      throw new ResponseStatusException(
          HttpStatus.BAD_REQUEST, "Livro nao encontrado");
    }
\end{lstlisting}
\end{frame}

\begin{frame}[fragile,shrink=5]{POST de vários --- gravar}
  \onde{Ainda no mesmo método}

\begin{lstlisting}[language=Java,basicstyle=\ttfamily\scriptsize]
    if (livro.getQuantidade() < item.getQuantidade()) {
      throw new ResponseStatusException(
          HttpStatus.BAD_REQUEST, "Quantidade insuficiente");
    }
  }

  for (int i = 0; i < pedido.getItens().size(); i++) {
    ItemEmprestimo item = pedido.getItens().get(i);
    Livro livro = acharLivro(item.getLivroId());
    livro.setQuantidade(
        livro.getQuantidade() - item.getQuantidade());
  }

  pedido.setId(Dados.proximoIdEmprestimo);
  Dados.proximoIdEmprestimo = Dados.proximoIdEmprestimo + 1;
  Dados.emprestimos.add(pedido);
  return pedido;
}
\end{lstlisting}
\end{frame}

\begin{frame}{Conferindo o passo 2}
  \begin{itemize}
    \item [ ] A tela Vários livros deixa acrescentar dois títulos
    \item [ ] Confirmar registra e as duas quantidades caem
    \item [ ] O histórico mostra os dois nomes
    \item [ ] O empréstimo de um livro ainda funciona
  \end{itemize}

  \vspace{0.5em}
  Se não der tempo agora, a gente faz a mesma lista\\
  no StockSales, depois do intervalo.
\end{frame}

\begin{frame}{Checklist da Parte 1}
  \begin{itemize}
    \item [ ] Estou no projeto \texttt{biblioteca}?
    \item [ ] \texttt{Dados} tem a lista de empréstimos?
    \item [ ] \texttt{/api/emprestimos} existe?
    \item [ ] A quantidade do livro diminui?
    \item [ ] A tela recusa quando não há exemplar?
  \end{itemize}
\end{frame}

\begin{frame}{Intervalo}
  \begin{center}
    \Large Retomamos às \textbf{20h20}\\[0.8em]
    \normalsize
    Fechem a \texttt{biblioteca}.\\
    Parem o Spring.\\
    Na segunda parte abrimos o \textbf{StockSales}.
  \end{center}
\end{frame}

% #############################################################################
\section{Parte 2 --- StockSales}

\begin{frame}{Mudamos de projeto}
  \onde{Agora é o projeto \textbf{StockSales}.\\
        A biblioteca já está encerrada.}

  Parem o Spring da biblioteca.\\
  Abram a pasta do StockSales e subam de novo.
\end{frame}

\begin{frame}{Onde o StockSales está}
  A loja já tem produto e cliente.

  \vspace{0.4em}
  \begin{itemize}
    \item Listas em \texttt{Dados}
    \item Telas de cadastro
    \item Telas de venda já prontas, com Axios
  \end{itemize}

  \vspace{0.5em}
  Ainda \textbf{não} tem a classe \texttt{Venda}.\\
  Ainda \textbf{não} tem o Controller de venda.

  \vspace{0.5em}
  É o empréstimo, com outro nome, e com preço.
\end{frame}

\begin{frame}{O desenho da venda de hoje}
  \begin{center}
  \begin{tikzpicture}[
      box/.style={draw, rounded corners, align=center,
                  minimum width=2.6cm, minimum height=0.9cm,
                  font=\small},
      seta/.style={-{Latex}, thick}
    ]
    \node[box, fill=blue!15] (cli) {Cliente};
    \node[box, fill=yellow!25, right=2.0cm of cli] (ven) {Venda};
    \node[box, fill=green!20, right=2.0cm of ven] (pro) {Produto};

    \draw[seta] (cli) -- (ven);
    \draw[seta] (pro) -- (ven);
  \end{tikzpicture}
  \end{center}

  \vspace{0.6em}
  Um cliente tem várias vendas.\\
  Um produto aparece em várias vendas.\\
  Cada venda, neste primeiro passo, aponta para \textbf{um} produto.
\end{frame}

\begin{frame}{O que a venda tem a mais}
  Igual ao empréstimo:
  \begin{itemize}
    \item \texttt{clienteId}
    \item \texttt{produtoId}
    \item \texttt{quantidade}
  \end{itemize}

  \vspace{0.4em}
  Extra da loja:
  \begin{itemize}
    \item \texttt{total} --- preço vezes quantidade
    \item o estoque do produto diminui
  \end{itemize}

  \vspace{0.4em}
  O JavaScript \textbf{não} calcula o total.\\
  Quem calcula é o servidor.
\end{frame}

\begin{frame}[fragile]{Vendas em Dados}
  \onde{Projeto \textbf{StockSales}\\
        Arquivo \texttt{Dados.java}, acrescentar}

\begin{lstlisting}[language=Java,basicstyle=\ttfamily\footnotesize]
public static ArrayList<Venda> vendas = new ArrayList<>();
public static int proximoIdVenda = 1;
\end{lstlisting}

  Import de \texttt{Venda} no topo.
\end{frame}

\begin{frame}[fragile]{Classe Venda}
  \onde{Arquivo novo:\\
        \texttt{src/main/java/com/entra21/stocksales/model/Venda.java}}

\begin{lstlisting}[language=Java,basicstyle=\ttfamily\footnotesize]
public class Venda {
  private int id;
  private int clienteId;
  private int produtoId;
  private int quantidade;
  private double total;

  public Venda() {
  }
}
\end{lstlisting}

  Construtor vazio e os gets e sets dos cinco atributos.
\end{frame}

\begin{frame}[fragile]{O JSON da nova venda}
  A página \texttt{vendas-nova.html} já manda:

\begin{lstlisting}[basicstyle=\ttfamily\footnotesize]
{"clienteId": 1, "produtoId": 2, "quantidade": 3}
\end{lstlisting}

  \vspace{0.4em}
  O servidor:
  \begin{enumerate}
    \item Acha cliente e produto
    \item Se \texttt{estoque < quantidade}, recusa
    \item \texttt{total = preco * quantidade}
    \item Diminui o estoque
    \item Guarda a venda
  \end{enumerate}
\end{frame}

\begin{frame}[fragile]{acharCliente e acharProduto}
  \onde{\texttt{controller/VendaController.java}, arquivo novo}

  Copiem a ideia da biblioteca.

\begin{lstlisting}[language=Java,basicstyle=\ttfamily\scriptsize]
public Produto acharProduto(int id) {
  for (int i = 0; i < Dados.produtos.size(); i++) {
    if (Dados.produtos.get(i).getId() == id) {
      return Dados.produtos.get(i);
    }
  }
  return null;
}
\end{lstlisting}

  A função \texttt{acharCliente} é a mesma, na lista
  \texttt{Dados.clientes}.
\end{frame}

\begin{frame}[fragile]{GET e começo do POST}
  \onde{Ainda em \texttt{VendaController.java}}

\begin{lstlisting}[language=Java,basicstyle=\ttfamily\scriptsize]
@GetMapping("/api/vendas")
public ArrayList<Venda> listar() {
  return Dados.vendas;
}

@PostMapping("/api/vendas")
public Venda cadastrar(@RequestBody Venda pedido) {
  Cliente cliente = acharCliente(pedido.getClienteId());
  if (cliente == null) {
    throw new ResponseStatusException(
        HttpStatus.BAD_REQUEST, "Cliente nao encontrado");
  }
\end{lstlisting}
\end{frame}

\begin{frame}[fragile,shrink=5]{POST --- estoque, total, gravação}
  \onde{Ainda no método \texttt{cadastrar}}

\begin{lstlisting}[language=Java,basicstyle=\ttfamily\scriptsize]
  Produto produto = acharProduto(pedido.getProdutoId());
  if (produto == null) {
    throw new ResponseStatusException(
        HttpStatus.BAD_REQUEST, "Produto nao encontrado");
  }

  if (produto.getEstoque() < pedido.getQuantidade()) {
    throw new ResponseStatusException(
        HttpStatus.BAD_REQUEST, "Estoque insuficiente");
  }

  pedido.setId(Dados.proximoIdVenda);
  Dados.proximoIdVenda = Dados.proximoIdVenda + 1;
  pedido.setTotal(produto.getPreco() * pedido.getQuantidade());
  produto.setEstoque(
      produto.getEstoque() - pedido.getQuantidade());
  Dados.vendas.add(pedido);
  return pedido;
}
\end{lstlisting}
\end{frame}

\begin{frame}{O caminho da primeira venda}
  \begin{enumerate}
    \item A pessoa escolhe Ana, Mouse Gamer, quantidade 2
    \item JS manda POST \texttt{/api/vendas}
    \item Servidor acha cliente e produto em \texttt{Dados}
    \item Confere estoque, calcula total, diminui estoque
    \item A tela de produtos, ao recarregar, já mostra
          o estoque menor
  \end{enumerate}
\end{frame}

\begin{frame}{Conferindo a venda de um produto}
  \begin{itemize}
    \item [ ] Nova venda baixa o estoque
    \item [ ] O histórico mostra o nome e o total
    \item [ ] Vender acima do estoque mostra a mensagem
    \item [ ] O total não vem digitado na tela
  \end{itemize}
\end{frame}

\begin{frame}{Venda com vários produtos}
  Mesma evolução da biblioteca.

  \vspace{0.5em}
  \begin{itemize}
    \item Classe \texttt{ItemVenda}: \texttt{produtoId},
          \texttt{quantidade}, \texttt{subtotal}
    \item Campo \texttt{itens} na \texttt{Venda}
    \item POST \texttt{/api/vendas/itens}
    \item A tela \texttt{vendas-varios.html} já existe
  \end{itemize}

  \vspace{0.5em}
  O POST de um produto permanece.
\end{frame}

\begin{frame}[fragile]{Classe ItemVenda}
  \onde{Arquivo novo: \texttt{model/ItemVenda.java}}

\begin{lstlisting}[language=Java,basicstyle=\ttfamily\footnotesize]
public class ItemVenda {
  private int produtoId;
  private int quantidade;
  private double subtotal;

  public ItemVenda() {
  }
}
\end{lstlisting}

  Gets e sets dos três atributos.\\
  O \texttt{subtotal} a tela não manda. O servidor preenche.
\end{frame}

\begin{frame}[fragile]{O JSON de vários produtos}
\begin{lstlisting}[basicstyle=\ttfamily\footnotesize]
{
  "clienteId": 1,
  "itens": [
    {"produtoId": 2, "quantidade": 1},
    {"produtoId": 4, "quantidade": 2}
  ]
}
\end{lstlisting}

  \vspace{0.4em}
  Primeiro \texttt{for}: conferir estoque de cada item.\\
  Segundo \texttt{for}: \texttt{subtotal}, somar o total, baixar estoque.
\end{frame}

\begin{frame}[fragile]{Somar o total}
  \onde{No segundo \texttt{for} de \texttt{cadastrarVarios}}

\begin{lstlisting}[language=Java,basicstyle=\ttfamily\scriptsize]
double total = 0;
for (int i = 0; i < pedido.getItens().size(); i++) {
  ItemVenda item = pedido.getItens().get(i);
  Produto produto = acharProduto(item.getProdutoId());
  item.setSubtotal(produto.getPreco() * item.getQuantidade());
  total = total + item.getSubtotal();
  produto.setEstoque(
      produto.getEstoque() - item.getQuantidade());
}
pedido.setTotal(total);
\end{lstlisting}

  \texttt{total = total + ...} vai juntando cada subtotal.
\end{frame}

\begin{frame}{Conferindo vários produtos}
  \begin{itemize}
    \item [ ] Ana leva Mouse e Headset na mesma venda
    \item [ ] Os dois estoques caem
    \item [ ] O total é a soma dos dois subtotais
    \item [ ] Se um produto não tem estoque, a venda inteira
          é recusada
    \item [ ] A venda de um produto continua funcionando
  \end{itemize}
\end{frame}

\begin{frame}{O que não entra hoje}
  \begin{itemize}
    \item Cancelar venda e devolver estoque
    \item Banco de dados
    \item Login
    \item Anotações JPA
  \end{itemize}

  \vspace{0.5em}
  \begin{block}{Meta no StockSales}
    Venda de um item, com estoque baixando e total no servidor.\\
    Se der tempo, vários itens na mesma venda.
  \end{block}
\end{frame}

\begin{frame}{Critérios de aceite}
  \begin{itemize}
    \item [ ] Venda grava \texttt{clienteId} e \texttt{produtoId}
    \item [ ] Total é calculado no servidor
    \item [ ] Estoque diminui
    \item [ ] Estoque insuficiente recusa a venda
    \item [ ] Histórico mostra o \textbf{nome}, não só o id
  \end{itemize}
\end{frame}

% #############################################################################
\section{Recapitulando}

\begin{frame}{O que vimos hoje}
  \begin{enumerate}
    \item 1{:}N guarda um id no lado dos muitos
    \item Quando dois cadastros se encontram, nasce uma classe nova
    \item Essa classe guarda os dois ids e a quantidade
    \item Uma lista de itens deixa vários produtos na mesma venda
    \item Conferir tudo antes de baixar estoque
  \end{enumerate}
\end{frame}

\begin{frame}{Os dois desenhos, lado a lado}
  \begin{center}
  \begin{tabular}{ll}
    \toprule
    biblioteca & StockSales \\
    \midrule
    Leitor e Livro & Cliente e Produto \\
    Emprestimo no meio & Venda no meio \\
    quantidade do livro cai & estoque do produto cai \\
    lista de itens & lista de itens, com total \\
    \bottomrule
  \end{tabular}
  \end{center}

  \vspace{0.6em}
  Mesma ideia. Nomes diferentes.
\end{frame}

\begin{frame}{Na próxima aula}
  \begin{block}{Aula 07}
    \begin{itemize}
      \item Validação dos dados que chegam
      \item Mensagem de erro mais clara
      \item Organizar o que já cresceu no StockSales
    \end{itemize}
  \end{block}

  \vspace{0.6em}
  Tragam o StockSales com a venda de um item funcionando.
\end{frame}

\begin{frame}
  \begin{center}
    \Huge Obrigado!\\[1.2em]
    \normalsize Mauricio Duailibi Neto\\
    Turma Java Entra21
  \end{center}
\end{frame}

\end{document}
